% SHARD-ID: AQM-53-MINIMAL-FERMION-GEOMETRY-CATALOGUE
% SHARD-TITLE: Minimal Fermion Geometry Catalogue
% SHARD-SUMMARY: Separates ordinary points, even nilpotents, odd variables, parity-shifted differentials, and spin data.
% SHARD-SUMMARY: Proves the minimal one-mode CAR, the dual-number differential calculation, and the finite-field scalar caveat.
% SHARD-SUMMARY: Gives concrete \(k=\mathbb F_3\) catalogue entries for the superpoint, the dual-number point, and a two-puncture log-spin curve.
% SHARD-KEYWORDS: fermion, superpoint, parity shift, Kahler differentials, spin structure, log form, CAR, F3

\section{Minimal Fermion Geometry Catalogue}
\label{sec:minimal-fermion-geometry-catalogue}

\subsection{Status and Source Anchors}
This shard is a definition catalogue for later fermion constructions.  It
does not assert that ordinary arithmetic quantum fields automatically contain
fermions.  The CAR input is AQM-19 and
\path{references/canonical_fields/SOURCES.md}.  Supergeometric odd variables,
reduced spaces, parity reversal, and Clifford quantization are sourced from
\path{references/supergeometry/SOURCES.md}.  K\"ahler differentials use
\path{references/algebraic_geometry/stacks_project_algebra.tex}, lines
33793--33853, 33856--33866, 34103--34120, 34310--34315, and 34392--34415.
Spin structures as square roots of the canonical line bundle use
\path{references/spin_geometry/SOURCES.md}.  Logarithmic differentials use
\path{references/log_geometry/SOURCES.md}.  The nilpotent Weyl layer is
AQM-36.  The identity \(\Gamma(\mathbf P^1,\mathcal O)=k\) is sourced by
\path{references/algebraic_geometry/SOURCES.md}.

\subsection{Lamport Dependency Skeleton}
\[
\begin{array}{rcl}
D1&:&\text{One complex fermion mode is the one-mode CAR representation.}\\
L1&:&D1\text{ satisfies }c^2=(c^*)^2=0,\ \{c,c^*\}=I.\\
D2&:&\text{A supercommutative }k\text{-algebra is }\mathbb Z/2\text{-graded}.\\
L2&:&\operatorname{char}k\ne2\text{ and }\theta\text{ odd imply }\theta^2=0.\\
D3&:&\mathsf{Spt}_k=(\Spec k,k[\theta])\text{ is the algebraic superpoint.}\\
D4&:&\Pi\Omega^1_{A/k}\text{ is the cotangent parity-shift mode source.}\\
L3&:&\Omega^1_{k/k}=0.\\
L4&:&\Omega^1_{k[\epsilon]/(\epsilon^2)/k}\cong k\,\mathrm d\epsilon
      \text{ when }\operatorname{char}k\ne2.\\
D5&:&(C,D,S,\mu:S^{\otimes2}\cong\Omega^1_C(\log D))
      \text{ is a log-spin datum.}\\
L5&:&(\mathbf P^1_k,\{0,\infty\},\mathcal O,\mathrm d\log x)
      \text{ is a log-spin datum.}\\
L6&:&\mathbb F_3\text{-vector spaces do not canonically scalar-extend to }
      \mathbb C\text{-Hilbert spaces.}\\
T1&:&\Spec k\text{ has no forced fermion in this catalogue.}\\
T2&:&\mathsf{Spt}_{\mathbb F_3},\ \Pi\Omega^1_{\mathbb F_3[\epsilon]/(\epsilon^2)/\mathbb F_3},
      \text{ and }(\mathbf P^1_{\mathbb F_3},\{0,\infty\},\mathcal O)
      \text{ give one mode each.}
\end{array}
\]

\subsection{The One-Mode CAR D1--L1}
Let
\[
  \mathcal H_1=\mathbb C|0\rangle\oplus\mathbb C|1\rangle .
\]
Define linear maps \(c,c^*:\mathcal H_1\to\mathcal H_1\) by
\[
  c|0\rangle=0,\quad c|1\rangle=|0\rangle,\quad
  c^*|0\rangle=|1\rangle,\quad c^*|1\rangle=0 .
\]
The catalogue phrase ``one fermion mode'' means this CAR representation, or a
unitarily isomorphic copy.  For \(r\) modes it means the antisymmetric Fock
space \(\bigwedge^\bullet\mathbb C^r\), as in AQM-19.

\paragraph{Lemma \(L1\).}
The operators \(c,c^*\) satisfy
\[
  c^2=(c^*)^2=0,\qquad cc^*+c^*c=I_{\mathcal H_1}.
\]

\emph{Proof.}
Both \(c^2\) and \((c^*)^2\) kill \(|0\rangle\) and \(|1\rangle\), so they are
zero.  Also
\[
  (cc^*+c^*c)|0\rangle=c|1\rangle+0=|0\rangle
\]
and
\[
  (cc^*+c^*c)|1\rangle=0+c^*|0\rangle=|1\rangle .
\]
The two basis vectors span \(\mathcal H_1\), hence \(cc^*+c^*c=I\).

\subsection{Odd Variables and the Superpoint D2--D3}
Assume \(\operatorname{char}k\ne2\) in this subsection.  A
supercommutative \(k\)-algebra is a \(\mathbb Z/2\)-graded algebra
\[
  A=A_{\bar0}\oplus A_{\bar1}
\]
such that homogeneous elements satisfy
\[
  uv=(-1)^{|u||v|}vu .
\]
Elements of \(A_{\bar0}\) are even; elements of \(A_{\bar1}\) are odd.

\paragraph{Lemma \(L2\).}
If \(\theta\in A_{\bar1}\), then \(\theta^2=0\).

\emph{Proof.}
Supercommutativity gives \(\theta\theta=(-1)\theta\theta\).  Hence
\(2\theta^2=0\).  Since \(2\in k^\times\), multiplication by \(2\) is
injective, so \(\theta^2=0\).

Define
\[
  k[\theta]=k\oplus k\theta,\qquad |\theta|=\bar1,\qquad \theta^2=0 .
\]
The algebraic superpoint is the locally ringed graded object
\[
  \mathsf{Spt}_k=(\Spec k,k[\theta]).
\]
Its reduced ordinary point is obtained by quotienting by the odd ideal:
\[
  k[\theta]/(\theta)\cong k .
\]
For \(k=\mathbb F_3\), the ungraded ring underlying \(k[\theta]\) is
isomorphic to the even dual-number ring \(k[\epsilon]/(\epsilon^2)\), by
\(\theta\mapsto\epsilon\).  The difference is the parity assignment, not the
ungraded multiplication table.

The CAR quantization convention for \(\mathsf{Spt}_k\) assigns the odd line
\(k\theta\) one complex fermion mode \(D1\).  This is a catalogue convention:
the standard Hilbert-space CAR is complex, while the odd geometric line is a
\(k\)-vector space.

\subsection{Cotangent Parity Shift D4--L4}
Let \(A\) be a commutative \(k\)-algebra.  The module of K\"ahler differentials
\(\Omega^1_{A/k}\) is the universal target for \(k\)-derivations
\(A\to M\), as recorded in the Stacks source.  The cotangent parity-shift
catalogue entry is
\[
  \Pi\Omega^1_{A/k},
\]
meaning that the module \(\Omega^1_{A/k}\) is treated as odd.  If
\(\dim_k\Omega^1_{A/k}=r<\infty\), the standard complex CAR realization uses
\(r\) complex fermion modes.

\paragraph{Lemma \(L3\).}
\(\Omega^1_{k/k}=0\).

\emph{Proof.}
Every \(k\)-derivation \(D:k\to M\) kills all elements of \(k\) by definition.
Thus the functor \(M\mapsto\operatorname{Der}_k(k,M)\) is the zero functor.
The zero module represents this functor, so the universal module
\(\Omega^1_{k/k}\) is zero.

\paragraph{Lemma \(L4\).}
Let \(A=k[\epsilon]/(\epsilon^2)\) and assume \(\operatorname{char}k\ne2\).
Then
\[
  \Omega^1_{A/k}\cong k\,\mathrm d\epsilon .
\]

\emph{Proof.}
Let \(S=k[t]\) and \(I=(t^2)\), so \(A=S/I\).  The Stacks polynomial-ring
lemma gives \(\Omega^1_{S/k}=S\,\mathrm dt\).  The quotient exact sequence for
differentials gives
\[
  I/I^2\longrightarrow \Omega^1_{S/k}\otimes_S A
  \longrightarrow \Omega^1_{A/k}\longrightarrow0,
\]
where \(t^2\mapsto \mathrm d(t^2)=2t\,\mathrm dt\).  Therefore
\[
  \Omega^1_{A/k}\cong A\,\mathrm d\epsilon/(2\epsilon\,\mathrm d\epsilon).
\]
Since \(2\in k^\times\), this is
\[
  A\,\mathrm d\epsilon/(\epsilon\,\mathrm d\epsilon)\cong (A/(\epsilon))\,
  \mathrm d\epsilon\cong k\,\mathrm d\epsilon .
\]

For \(k=\mathbb F_3\), the even thickened point
\(\Spec(\mathbb F_3[\epsilon]/(\epsilon^2))\) has two distinct catalogue
layers:
\[
  \ell^2(\mathbb F_3[\epsilon]/(\epsilon^2))
\]
is the \(9\)-dimensional even Artin-Weyl layer of AQM-36, while
\(\Pi\Omega^1_{A/k}\cong\mathbb F_3\,\mathrm d\epsilon\) supplies one optional
fermion mode by the cotangent parity-shift convention.

\subsection{Log-Spin Datum D5--L5}
Let \(C\) be a smooth curve over \(k\), and let \(D\) be an effective divisor
for which the logarithmic one-forms \(\Omega^1_C(\log D)\) are defined.  A
log-spin datum in this catalogue is
\[
  (C,D,S,\mu),\qquad
  \mu:S^{\otimes2}\xrightarrow{\sim}\Omega^1_C(\log D),
\]
where \(S\) is a line bundle on \(C\).  Its algebraic mode source is
\[
  \Gamma(C,S).
\]
If \(\dim_k\Gamma(C,S)=r<\infty\), the standard complex CAR realization uses
\(r\) complex fermion modes.  Atiyah's source supports the square-root spin
dictionary over compact complex curves; the logarithmic variant is local.

\paragraph{Lemma \(L5\).}
Let \(C=\mathbf P^1_k\) with affine coordinate \(x\), and let
\[
  D=\{0,\infty\}.
\]
Then \(\Omega^1_C(\log D)\cong\mathcal O_C\), generated by
\[
  \omega=\mathrm d\log x .
\]
Consequently \((C,D,\mathcal O_C,\mu)\), with
\(\mu(1\otimes1)=\omega\), is a log-spin datum.

\emph{Proof.}
On \(U_0=\Spec k[x]\), the divisor point \(0\) has local equation \(x\).
The Stacks log source says logarithmic one-forms are generated by ordinary
forms and \(\mathrm d\log x=x^{-1}\mathrm dx\).  Since
\(\mathrm dx=x\,\mathrm d\log x\), the module is generated by
\(\mathrm d\log x\) on \(U_0\).  On \(U_\infty=\Spec k[u]\), with \(u=x^{-1}\),
the divisor point \(\infty\) has local equation \(u\), and
\[
  \mathrm d\log x=\mathrm d\log(u^{-1})=-\mathrm d\log u .
\]
Thus the same form generates the logarithmic module on \(U_\infty\).  The two
descriptions agree on \(U_0\cap U_\infty\), so they glue to a global basis of
\(\Omega^1_C(\log D)\).  Therefore \(\Omega^1_C(\log D)\cong\mathcal O_C\).
The map \(\mu:\mathcal O_C^{\otimes2}\to\Omega^1_C(\log D)\) sending
\(1\otimes1\) to \(\omega\) is an isomorphism.

\subsection{Finite-Field Scalar Caveat L6}
\paragraph{Lemma \(L6\).}
There is no unital field homomorphism \(\mathbb F_3\to\mathbb C\).

\emph{Proof.}
If \(\iota:\mathbb F_3\to\mathbb C\) were unital, then
\[
  0=\iota(0)=\iota(1+1+1)=1+1+1=3,
\]
which is false in \(\mathbb C\).

Therefore an \(\mathbb F_3\)-vector-space mode source \(V\) does not
canonically become a complex one-particle Hilbert space by scalar extension.
The catalogue records \(r=\dim_{\mathbb F_3}V\) and then uses a chosen complex
Hilbert space \(\mathbb C^r\) for the standard CAR realization.  A genuinely
\(\mathbb F_3\)-linear Clifford theory is a different algebraic layer and is
not defined here.

\subsection{Catalogue Consequences T1--T2}
\paragraph{Theorem \(T1\).}
The ordinary point \(\Spec k\) has no forced fermion mode in this catalogue.

\emph{Proof.}
The superpoint entry requires an odd structure sheaf \(k[\theta]\); the
ordinary structure sheaf \(k\) has no odd part.  The cotangent parity-shift
entry gives no mode by Lemma \(L3\).  The log-spin entry requires a curve,
a divisor, and a square-root line bundle; none is part of the ordinary point
\(\Spec k\).  Thus no catalogue mechanism above assigns a nonzero fermion
mode to \(\Spec k\) without extra data.

\paragraph{Theorem \(T2\).}
For \(k=\mathbb F_3\), the following are the minimal one-mode entries:
\[
\begin{array}{c|c|c|c}
\text{geometric datum}&\text{mode source}&\mathbb F_3\text{-dimension}
&\text{complex CAR Hilbert space}\\
\hline
\mathsf{Spt}_{\mathbb F_3}&\mathbb F_3\theta&1&\mathcal H_1\\
\Spec(\mathbb F_3[\epsilon]/(\epsilon^2))\text{ with }\Pi\Omega^1
&\mathbb F_3\,\mathrm d\epsilon&1&\mathcal H_1\\
(\mathbf P^1_{\mathbb F_3},\{0,\infty\},\mathcal O)&
\Gamma(\mathbf P^1,\mathcal O)=\mathbb F_3&1&\mathcal H_1
\end{array}
\]
The ordinary closed point \(\Spec\mathbb F_3\) contributes zero fermion modes.

\emph{Proof.}
The superpoint row is Definition \(D3\).  The dual-number row is Lemma \(L4\).
The log-spin row is Lemma \(L5\) plus the cited projective-space cohomology
calculation \(\Gamma(\mathbf P^1,\mathcal O)=\mathbb F_3\).  Lemma \(L6\)
explains why each row determines the standard complex CAR Hilbert space only
through the dimension-one mode count.  The last sentence is Theorem \(T1\)
with \(k=\mathbb F_3\).
